\begin{topic}[id=ros2_micro,
             difficulty={Anspruchsvoll},
             type={Systemanalyse + Middleware},
             prerequisites={Grundlagen Robotik oder eingebettete Systeme; Basiswissen Linux und verteilte Systeme hilfreich},
             practice={optional},
             owner={Dozententeam},
             updated={2026-04-09},
             tags={ ROS 2, micro-ROS, DDS, QoS, RT, Executor, Sensorik, Aktorik }]{ROS 2 \& micro-ROS}

\TopicDescription{ROS 2 ist der Standard-Stack für Robotik; micro-ROS bringt ROS-Themen auf Microcontroller. Du lernst die Architektur (Nodes, Topics, Services, Actions), DDS/QoS als Transportbasis und die Besonderheiten echtzeitnaher Ausführung (Executors, RMW). Wir beleuchten, wann micro-ROS sinnvoll ist, wie es sich mit ROS 2 integriert und welche Ressourcenfallen (Speicher, Timing) realistisch auftreten. Ziel ist eine klare Vorstellung, wie man ROS 2-Graphen robust aufsetzt und QoS für Sensorik/Steuerung auswählt.}

\TopicLeitfragen{\item Welche QoS-Profile sind für Sensor- vs. Steuerpfade geeignet?
\item Wann lohnt micro-ROS statt „voller“ ROS 2?
\item Wie strukturiert man stabile Launch/Composability?
}
\TopicScope{\item Kein Gazebo/Simulations-Deep-Dive.
\item Keine Treiberentwicklung im Detail.
}
\TopicPractice{Optionale Praxisidee: Kleiner ROS 2-Graph mit zwei QoS-Varianten (Latency/Drop vergleichen).}

\TopicKeywords{ROS 2, micro-ROS, DDS, QoS, RT, Executor, Sensorik, Aktorik}

\nocite{ros2Docs,microRosDocs,ros2Design}
\end{topic}
